%!TEX root=../mythesis.tex
% Chapter Template

\chapter{Literature Review} % Main chapter title
\chaptermark{Literature Review}  % replace the chapter name with its abbreviated form
\label{ch:literature_review} % Change X to a consecutive number; for referencing this chapter elsewhere, use \ref{ChapterX}
In this chapter, we provide an overview of the related works on resource allocation and job scheduling in single-cluster and distributed environments. In this dissertation, we focus on developing algorithms that achieve better performance in fairness and efficiency with the utilization of the network resource and trade-off between these two objectives without the network. Therefore, we review some literature based on different fairness and efficiency objectives as well as different problem scenarios.


\section{Fairness}
% To be modified further.
In modern computing-intensive systems, fairness is a critical performance metric that ensures equitable resource allocation among users or jobs. Fairness can be defined in various ways, depending on the specific context and objectives of the system. In this section, we review the existing literature on fairness in both single-cluster and distributed job execution scenarios.

\subsection{Fairness in Single-cluster System}
A multitude of notions exist for different types of fair resource allocation. The most widely used resource allocation policy in the single-cluster system is the Proportional Sharing (PS) policy, which allocates resources (i.e., system processing capacities) evenly to jobs pending for processing. If the system is not aware of the specific workload but only knows the resource demand of each job, the PS policy degenerates to the Max-min  Fairness (MMF) allocation. Radunovic \textit{et al.} \cite{radunovic2007unified} proves that the max-min fair allocation exists uniquely at any compact and convex set and presents general algorithms for computing the MMF resource allocation in the single-cluster system, which aims to maximize the most unsatisfied job using a water-filling approach. 

In addition, there are other notions of fairness such as the proportional fairness (allocating resources proportionally to the jobs' demands) and weighted max-min fairness. Mo \textit{et al.} \cite{mo2000fair} presented $\alpha$-Fairness to unify various existing notions of fairness. Suppose the resource allocation is a vector $\mathbf{x} = (x_1, x_2, \dots, x_n)$ where each element $x_i$ represents the resource allocation to the $i$-th user, $\alpha$-Fairness defines the utility function as $U_\alpha (\mathbf{x}) = \frac{\mathbf{x}^{1-\alpha}}{1-\alpha}$ if $\alpha \ne 1$ and $U_\alpha (\mathbf{x}) = \ln \mathbf{x}$ if $\alpha = 1$. If $\alpha=0$, the $\alpha$-Fairness becomes efficiency-targeted which tries to maximize the total system utility; when $\alpha=1$, it becomes the proportional fairness; when $\alpha \rightarrow \infty$, it becomes the max-min fairness. In addition, Ghodsi \textit{et al.} \cite{choosy} addressed task placement constraints (i.e., each task can only be assigned to one of the certain subset of all machines in the system) and presented Constrained Max-min Fairness (CMMF), which tries to meet the fairness target by recursively satisfying the users based on the order of dissatisfaction. 



\subsection{Fairness for Multiple Types of Resources}
There are also studies working on the multi-resource allocation case. Cobb Ghodsi \textit{et al.} \cite{DRF} proposed Dominant Resource Fairness (DRF) which decides a user's resource allocation by its dominant share, i.e., the maximum proportional share that the user demands among all resource types. Zahedi and Lee \cite{REF} solved the multi-resource allocation problem by proportionally allocating different resources based on the re-scaled Cobb-Douglas Utility and named it Resource Elasticity Fairness (REF). REF defines each user's utility as $u_i(x_i) = \alpha_{i0} \prod_{r=1}^Rx_{ir}^{a_{ir}}$, where $x_{ir}$ denotes the allocation of resource type $r$ to $i$-th user and $\alpha_{ir}$ denotes the preference of the $i$-th user to resource type $r$ ($\alpha_{i0}$ represents the overall preference). Gutman and Nisan \cite{gutman2012GDF} further summarized that DRF actually uses $L_{\inf}$ norms (i.e., the dominant share) to measure the utility of each user and presented Generalized Resource Fairness (GRF) which allows measuring the fairness of multi-resource allocation with a wider range of norms. Fossati \textit{et al.} \cite{fossati2022distributed} presents a distributed slice multi-resource allocation algorithm specialized for centralized 5G systems. Li \textit{et al.} \cite{li2025fair} studies the fair resource allocation based on various SLAs of users while improve the execution efficiency. Guo \textit{et al.} \cite{guo2024dynamic} presents a fair multi-resource allocation mechanism in cloud computing system with elastic user demands. Multi-resource allocation holds wide applications in cloud servers \cite{wang2014dominant}, heterogeneous GPU clusters \cite{mahajan2020themis}, 5G networks \cite{fossati2020multi} and LLM inference \cite{sheng2024fairness}. In particular, Friedman \textit{et al.} \cite{friedman2003fairness, friedman2003protective} presents the notion of protective scheduling, which sets the criterion of fairness to be that no job would finish later than its completion time under the PS policy. Under this fairness guarantee, they present a series of algorithms to further improve the job processing efficiency. %Our work also considers multiple types of resources, but we focus on utilizing network resources to enhance the fairness of compute resource allocation.
%Existing studies on fair resource allocation, such as Dominant Resource Fairness \cite{DRF} and Aggregate Max-min Fairness \cite{AMF}, either predominantly operate under single-site assumptions where all systemic resources are accessible to all jobs, or do not consider the potential of task transmission. Naive extension of single-site fairness at the localized level (i.e., applying the fairness principle at each site independently) is inherently flawed in distributed job execution case; for example, a resource-intensive job constrained to a low-capacity site may suffer from highly unfair allocations despite monopolizing local resources. In addition, transferring tasks (resource demands) in distributed job execution case can potentially improve fairness by allowing starving jobs to access resources from other sites, but it requires sophisticated mechanisms design to avoid excessive transmissions.



\section{Fairness and Efficiency in Distributed Job Execution Cases}
The task (job) scheduling problem in geo-distributed environments has been widely studied from a variety of system scenarios and performance objectives \cite{gautam2015survey, ghomi2017load, hedayati2023mapreduce, wu2025task}. Since the job load and the system capacity are distributed (mostly unevenly) across multiple sites, the problem becomes significantly more complex due to the additional dimension of site-specific resource constraints and inter-site communication overheads. Although the SRPT policy can be extended to such distributed job execution cases, it admits multiple algorithmic variations and they all no longer guarantees optimality. Hung \textit{et al.} \cite{hung2015scheduling} presents SWAG, a workload-aware strategy which schedules jobs based on their estimated completion times under given task placements, but it is efficiency-targeted and neglects to utilize network resources to move tasks. Guan \textit{et al.} \cite{guan2022task}, Zhao \textit{et al.} \cite{zhao2025data} study the task assignment problem given that the input data of each task can be replicated across multiple geo-distributed data centers, and they do not consider using network resources either. Hu \textit{et al.} \cite{hu2016flutter} and Pu \textit{et al.} \cite{pu2015low} study algorithms considering task transmissions, but they focus on improving the efficiency of a single job and do not study the scheduling strategy across jobs. Hung \textit{et al.} \cite{hung2018wide} proposes Tetrium, an entire system that incorporates both compute and network resources to improve the efficiency of map-reduce job execution, but it focuses on the task assignment of each single job and simply applies SRPT to schedule multiple jobs. Convolbo \textit{et al.} \cite{GEODIS} and Li \textit{et al.} \cite{li2020load} coordinates scheduling with transmission but focuses on minimizing the makespan of the entire system instead of the average response time of jobs. Li \textit{et al.} \cite{li2021mapreduce} improves the MapReduce task scheduling efficiency by monitoring the idle containers in the geo-distributed system and solving the task assignment problem using the classical Hungarian algorithm. Some research focuses on utilizing the game theory to minimize the system entity cost by assuming that users in the system have different interests, and thus they form a game model \cite{chen2014decentralized, hogade2022energy}. Recent studies have also adopted deep reinforcement learning to search for effective resource allocation strategies in fields like wireless networks \cite{wang2021incorporating}, edge computing \cite{li2022drl, mohajer2024dynamic}, satellite networks \cite{liao2020distributed, li2022multi} as well as conventional distributed data centers \cite{wei2022multi}. However, learning-based approaches have mainly focused on improving the efficiency of resource allocation (such as system utility, revenue and profit) or achieving load balancing instead of achieving the fairness of resource allocation.

Many efforts aim to minimize the bandwidth usage or the transmission time/cost in the system. For example, Hsieh \textit{et al.} \cite{hsieh2017gaia} presents Gaia to minimize the transmission cost between machine learning clusters while maintaining the model's performance, Chen \textit{et al.} \cite{chen2021optimizing} focuses on optimizing the transmission time of data along the given paths, and Shi \textit{et al.} \cite{shi2024adaptive} aims to save the network bandwidth usage on both intra-DC and inter-DC links while ensuring the QoS of microservices. In addition, Zhao \textit{et al.} \cite{zhao2019optimizing} models the scheduling problem in geo-distributed systems as an extension to Hypergraph Partition-based Scheduling which aims to save the WAN data transmission amount and improves system makespan, namely HPS+.  

In the context of fairness, Guan \textit{et al.} \cite{AMF} introduces a task assignment algorithm aimed at achieving max-min fairness in aggregate resource allocation among jobs in the distributed job execution. Similarly, Chen \textit{et al.} \cite{chen2018scheduling} investigates achieving max-min fairness in terms of the job response times in geo-distributed systems by task transmission, but imposes strict constraints that the number of tasks is no more than the number of computing slots in the system. Meskar \textit{et al.} \cite{meskar2022fair} presents a fair multi-resource allocation policy specialized for heterogenous servers in edge computing. Li \textit{et al.} studies the similar problem with access constraint on each user to the servers \cite{li2024fair}. In addition, there have also been researches investigating the job trace properties to design specialized scheduling or resource allocation algorithms \cite{tian2019characterizing, luo2021characterizing}.

\section{Efficiency}

\subsection{Shortest Remaining Processing Time (SRPT) and Others}
In the single-cluster scenario, the Shortest Remaining Processing Time (SRPT) policy is universally recognized as the optimal policy for maximizing job processing efficiency.
%Different from the works above, we focus on the dual objective of efficiency and long-term fairness in geo-distributed job execution without transmission.

%In geo-distributed systems, a job may have tasks spread over different sites. Hung \textit{et al.} \cite{hung2015scheduling} presented a scheduling algorithm for distributed job execution, which selects the job with the least overall demand to run first in order to minimize the average response time of jobs. This allocation policy is straightforward but apparently unfair, and could cause starvation for large-demand jobs. Different from the above work, we aim to achieve max-min fairness in the instantaneous resource allocation among jobs across geo-distributed sites given the availability of network resources. 